% !TeX root = ../main.tex
\section{บทนำ}
การลาหยุดเรียนของนิสิตต้องผ่านการพิจารณาจากผู้รับผิดชอบหลายฝ่าย
และเกี่ยวข้องกับการเข้าเรียนในแต่ละรายวิชา
ปัญหาที่พบคือการติดตามสถานะคำขอลาทำได้ยาก
และไม่มีการรวบรวมประวัติการลาไว้ในแหล่งเดียว
ส่งผลให้นิสิตไม่สะดวกในการตรวจสอบความคืบหน้าของคำขอ
ขณะที่อาจารย์ต้องค้นหาและรวบรวมข้อมูลจากเอกสารหลายส่วนเมื่อต้องการตรวจสอบการลาย้อนหลังหรือการลาในรายวิชาที่รับผิดชอบ

ปัจจุบันการดำเนินการลาของนิสิตใช้เอกสารกระดาษ
โดยนิสิตต้องจัดเตรียมใบลา แนบหลักฐาน
และดำเนินการส่งเอกสารให้ผู้รับผิดชอบพิจารณาตามขั้นตอนด้วยตนเอง
การติดตามผลจึงต้องติดต่อผู้เกี่ยวข้องเป็นรายครั้ง
และประวัติการลายังกระจายอยู่ตามเอกสารที่จัดเก็บ \cite{proposal}
นอกจากนี้ การตรวจสอบว่านิสิตที่ขาดเรียนในรายวิชามีการลาหรือไม่
ต้องนำข้อมูลจากใบลามาเปรียบเทียบกับใบเช็กชื่อ
จึงมีความจำเป็นต้องรวบรวมข้อมูลคำขอ สถานะการพิจารณา
และรายวิชาที่เกี่ยวข้องให้สามารถตรวจสอบร่วมกันได้

โครงงานนี้จึงพัฒนาเว็บแอปพลิเคชันระบบลานิสิตสำหรับคณะวิศวกรรมศาสตร์
กำแพงแสน มหาวิทยาลัยเกษตรศาสตร์
เพื่ออำนวยความสะดวกให้นิสิตและอาจารย์ในการดำเนินการและติดตามผลการลา
ระบบรองรับการยื่นคำขอและแนบหลักฐานผ่านเว็บไซต์
การพิจารณาตามลำดับการอนุมัติ
และการตรวจสอบสถานะพร้อมประวัติการลาในระบบเดียว
อีกทั้งให้อาจารย์ผู้สอนตรวจสอบข้อมูลการลาที่ได้รับอนุมัติในรายวิชาที่รับผิดชอบ
เพื่อนำไปเปรียบเทียบกับใบเช็กชื่อ
โดยกำหนดสิทธิ์การเข้าถึงข้อมูลตามบทบาทและความเกี่ยวข้องกับนิสิต

บทความนี้นำเสนอการออกแบบและพัฒนาระบบที่เชื่อมโยงกระบวนการยื่นคำขอ
การพิจารณา และการติดตามผลเข้ากับข้อมูลรายวิชา ผู้สอน
และอาจารย์ที่ปรึกษา
พร้อมแยกคลังข้อมูลรายงานออกจากฐานข้อมูลธุรกรรมเพื่อสนับสนุนการสรุปข้อมูลการลาตามบทบาทของผู้ใช้
เนื้อหาครอบคลุมแนวคิดและงานที่เกี่ยวข้อง วิธีการออกแบบระบบ
ผลการตรวจองค์ประกอบการทำงานและแนวทางประเมิน
ตลอดจนข้อเสนอแนะสำหรับการพัฒนาต่อไป
โดยผลที่นำเสนออาศัยการตรวจซอร์สโค้ดของระบบ \cite{repo}
ส่วนการประเมินประสิทธิภาพและผลการใช้งานกับผู้ใช้จริงยังต้องดำเนินการเพิ่มเติม
